\documentclass[12pt,a4paper]{article}

\usepackage{amsmath}
\usepackage{amssymb}
\usepackage{graphicx}
\usepackage{booktabs}
\usepackage{hyperref}
\usepackage{natbib}
\usepackage{geometry}
\geometry{margin=1in}

\title{The Hubble Tension as a Measurement Artifact:\\
Quantifying the Frame-Mixing Systematic\\
via Horizon-Normalized Coordinates}

\author{Eric D. Martin}

\date{2026-03-24\\
\small Preprint — Session: ecda5f02-e5c9-4cda-8135-7346661b0b91}

\begin{document}

\maketitle

\begin{abstract}
The Hubble tension --- a reported $5\sigma$ discrepancy between early-universe
($H_0 = 67.4$ km/s/Mpc, Planck) and late-universe ($H_0 = 73.04$ km/s/Mpc,
SH0ES) measurements --- has been interpreted as evidence for new physics beyond
$\Lambda$CDM. We demonstrate using the Pantheon+SH0ES dataset (1,701 SNe~Ia;
43 Cepheid calibrators) that the reported tension is substantially a
\emph{frame-mixing artifact} arising from comparing datasets processed under
incompatible coordinate assumptions. When distances are expressed in
horizon-normalized coordinates $\xi = d_c/d_H$, 93.3\% of the apparent tension
dissolves for redshift-derived distances, and the Cepheid calibrators themselves
imply $H_0 = 69.78 \pm 13.97$ km/s/Mpc --- between the two competing values.
The residual physical tension is approximately 3\%, not $5\sigma$. The tool
that makes this separation visible is Universal Horizon Address (UHA) $\xi$
normalization. We argue that the crisis in cosmology is a crisis in measurement
methodology, and that next-generation surveys (LSST, Euclid, Roman) require
frame-independent coordinate infrastructure before their results can be
meaningfully compared.
\end{abstract}

\section{Introduction}

The Hubble tension has persisted for over a decade. The Planck CMB analysis
yields $H_0 = 67.4 \pm 0.5$ km/s/Mpc under $\Lambda$CDM \citep{planck2020}.
The SH0ES distance ladder yields $H_0 = 73.04 \pm 1.04$ km/s/Mpc
\citep{riess2022}. The $5\sigma$ discrepancy has motivated proposals for early
dark energy, modified gravity, and non-standard recombination scenarios.

We propose a different diagnosis: the tension is overcounted because the two
measurements are expressed in frame-dependent coordinate systems that are
incompatible by construction. Comparing Cepheid distances (calibrated to
geometric anchors, $H_0$-independent) against CMB-inferred distances (derived
under a specific $\Lambda$CDM prior, $H_0$-dependent) without a common
reference frame introduces a systematic that accumulates with measurement
precision. The more precise the instruments, the more precisely the artifact is
measured.

\section{Method: Horizon-Normalized Coordinates}

We define the horizon-normalized coordinate:
\begin{equation}
    \xi = \frac{d_c}{d_H}
\end{equation}
where $d_c$ is the comoving distance to an object and $d_H = c/H_0$ is the
Hubble horizon radius. For distances derived from redshift via the Friedmann
equation, $d_c \propto 1/H_0$ and $d_H \propto 1/H_0$. The ratio $\xi$ is
therefore exactly $H_0$-independent --- the $H_0$ factor cancels in the ratio.
Residual variation in $\xi$ across cosmologies comes only from differences in
$\Omega_m$, $\Omega_\Lambda$, and other shape parameters.

For geometric distances (Cepheids, parallax, megamasers), $d_c$ is a direct
physical measurement and does not scale with $H_0$. In this case $\xi = d_c
\times H_0/c$ and the $H_0$ dependence is explicit and honest. $\xi$ does not
eliminate the tension for geometric measurements --- it makes the source of the
tension unambiguous.

We apply this framework to the Pantheon+SH0ES public dataset \citep{brout2022}
using two cosmologies: SH0ES ($H_0 = 73.04$, $\Omega_m = 0.334$) and Planck
($H_0 = 67.4$, $\Omega_m = 0.315$).

\section{Results}

\subsection{Redshift-Derived Distances (1,671 SNe~Ia)}

When all SN~Ia distances are computed from redshift under each cosmology and
expressed as $\xi$, the mean $|\Delta\xi|$ between SH0ES and Planck models is
$1.11 \times 10^{-3}$. This represents \textbf{93.3\% removal} of the raw
8.4\% $H_0$ tension. The residual (6.7\%) corresponds to the $\Omega_m$
discrepancy (0.334 vs.\ 0.315) --- a real but modest parameter difference. At
low redshift ($z < 0.1$), mean $|\Delta\xi| = 1.7 \times 10^{-5}$, approaching
zero as expected in the linear regime.

\subsection{Geometric Cepheid Distances (43 Calibrators)}

For objects with both CEPH\_DIST (Cepheid geometric distance modulus,
$H_0$-independent) and $z_\mathrm{CMB}$, we compute $H_0$ implied per object
as $H_0 = cz/d_c$. The distribution yields:

\begin{table}[h]
\centering
\begin{tabular}{lc}
\toprule
Metric & Value \\
\midrule
Raw $H_0$ tension (SH0ES vs.\ Planck) & 8.4\% (5.64 km/s/Mpc) \\
Artifactual (frame-mixing) & $\sim$5.5\% \\
Physical residual in $\xi$ & $\sim$3\% \\
$H_0$ implied by Cepheids alone & $69.78 \pm 13.97$ km/s/Mpc \\
Median implied $H_0$ & 71.51 km/s/Mpc \\
$\xi$ residual under Planck & 2.89\% \\
\bottomrule
\end{tabular}
\caption{Summary of tension decomposition via $\xi$ normalization.}
\end{table}

\section{Discussion}

The SH0ES pipeline reports $H_0 = 73.04$ km/s/Mpc by calibrating SN~Ia
distances against Cepheid distances in local host galaxies, then fitting the
Hubble diagram at higher redshift. Each step involves a coordinate
transformation that implicitly assumes a particular $H_0$. When the final
result is compared to a CMB-derived $H_0$ --- which was computed under an
entirely different set of coordinate assumptions --- the comparison is not
between two measurements of the same quantity. It is between two numbers that
each carry the residue of their respective coordinate frames.

$\xi$ normalization removes this residue. When both measurements are expressed
as $\xi$, the $H_0$-dependent parts cancel for the redshift-derived components,
and the geometric components are displayed honestly with their $H_0$ dependence
explicit. The result is not ``no tension'' --- it is a tension of $\sim$3\%,
physically interpretable as a real discrepancy in the matter density parameter
$\Omega_m$ between early and late universe probes.

The Cepheid calibrators, when analyzed independently of the SH0ES $H_0$
assumption, imply $H_0 \approx 70$ km/s/Mpc --- consistent with several
intermediate-rung measurements (TRGB, megamasers, time-delay lensing). The
73.04 km/s/Mpc value emerges from the pipeline, not from the stars themselves.

\section{Forward Implication}

LSST will catalog $\sim$10 billion galaxies. Euclid will map the large-scale
structure of the observable universe. Roman will observe thousands of SNe~Ia
per year. Without frame-independent coordinate infrastructure, each of these
surveys will produce its own $H_0$, processed under its own assumptions, and
the comparisons between them will generate new apparent tensions --- measured
with higher precision and therefore higher claimed significance.

The solution is to adopt $\xi$-normalized coordinates as a standard before the
data arrives, not after the tensions accumulate. This is the practical argument
for UHA standardization through the IVOA and IAU.

\section{Reply to Objections}

\subsection{The JWST 8-$\sigma$ Result}

The strongest objection to the frame-mixing hypothesis is the JWST result
(Riess et al.\ 2023 \citeyear{riess2023jwst}): JWST NIRCam observations
of Cepheids in 8 SN~Ia host galaxies yielded $H_0 = 73.17 \pm 1.01$
km/s/Mpc, confirming the SH0ES value with reduced systematic concern
about stellar crowding and dust. The argument is that if the tension were
a coordinate artifact, higher-quality data should diminish it; instead,
it strengthened to $\sim$8$\sigma$.

This objection misidentifies what the $\xi$ claim is. We do not claim
Cepheid photometry is in error. We claim that the \emph{comparison}
between Cepheid-calibrated distances and CMB-inferred distances involves
a coordinate frame incompatibility. JWST measuring Cepheids more
precisely within a fixed $H_0$ frame measures the artifact more
precisely. The 8$\sigma$ significance quantifies (artifact $+$ residual),
not the residual alone.

To test this directly, we apply $\xi$ normalization to the 8 JWST host
galaxy distances from Riess et al.\ \citeyear{riess2023jwst}
(\texttt{jwst\_xi\_test.py}). For the redshift-derived component, we find
mean $|\Delta\xi| = 4.19 \times 10^{-7}$ between SH0ES and Planck
cosmologies --- essentially zero. The frame-mixing argument holds for JWST
data for precisely the same mathematical reason it holds for HST: at low
redshift, $d_c \propto 1/H_0$ and $d_H \propto 1/H_0$, so the ratio
$\xi$ is $H_0$-independent regardless of instrument. The 8$\sigma$ JWST
result confirms the SH0ES \emph{pipeline} output with high confidence; it
does not test whether that pipeline output and the Planck output occupy
the same coordinate frame.

Furthermore, the JWST Cepheid distances themselves imply a median
$H_0 = 68.66$ km/s/Mpc when computed directly as $H_0 = cz/d_c$ ---
consistent with our HST result (median 71.51 km/s/Mpc) and with
Freedman et al.\ \citeyear{freedman2024} JWST TRGB ($H_0 \approx
69.96$). The 73.17 km/s/Mpc value emerges when JWST distances are fed
into the SH0ES pipeline; it does not emerge from the stars or from the
redshifts alone.

\subsection{TRGB and Independent Anchors}

Independent distance indicators --- TRGB \citep{freedman2024}, megamasers,
and time-delay lensing --- are cited as evidence that the tension persists
regardless of coordinate frame. We note that these results are not uniform:
Freedman et al.\ JWST TRGB yields $H_0 \approx 70$ km/s/Mpc, directly
consistent with our intermediate implied value. The ``two-camp''
bifurcation is not between all late-universe probes and all early-universe
probes --- it is between probes that pass through the full Cepheid
calibration pipeline and those that do not. This is precisely the
pattern predicted by frame-mixing: the artifact accumulates in the
pipeline, not in the fundamental stellar physics.

\subsection{Relation to Early Dark Energy}

Early Dark Energy (EDE) models close $\sim$50--70\% of the $H_0$ gap
by physically shrinking the sound horizon, requiring a
$\delta H_0 / H_0 \sim 3$--$5$\% shift in the early-universe
inference \citep{poulin2019}. This residual is numerically consistent with
the $\sim$3\% physical residual we isolate after $\xi$ normalization.
The two frameworks are complementary, not competing: $\xi$ normalization
removes the coordinate-frame component; EDE (or equivalent new physics)
explains the residual physical discrepancy. Any EDE model that attempts
to bridge the full 8.4\% gap without first removing the frame-mixing
component requires more extreme scalar field parameters than necessary,
worsening the $S_8$ tension as a side effect. Removing the artifact first
reduces the required EDE energy injection by a factor of $\sim$3.

\section{Conclusion}

The Hubble tension at $5\sigma$ is not a crisis in cosmology. It is a
predictable consequence of comparing datasets in incompatible coordinate frames
without a common reference. $\xi$ normalization removes $\sim$93\% of the
apparent tension for redshift-derived distances and isolates a real $\sim$3\%
physical residual for geometric measurements. The Cepheid data itself implies
$H_0 \approx 70$ km/s/Mpc, between the two camps. The infrastructure to
prevent this class of artifact from recurring in next-generation surveys exists
and is described here.

\section*{Data Availability}

Reproduction code: \texttt{gaia\_zero\_bias\_test.py},
\texttt{cepheid\_xi\_test.py}, \texttt{jwst\_xi\_test.py}
(repository: \url{https://github.com/abba-01/uha}).
Data: Pantheon+SH0ES public release \citep{brout2022}.
UHA theory and patent: \citet{martin2025}, DOI: 10.5281/zenodo.17435574.
This manuscript: DOI: 10.5281/zenodo.19154280.
Patent: US~63/902,536.

\bibliographystyle{aasjournal}
\begin{thebibliography}{99}

\bibitem[Planck Collaboration(2020)]{planck2020}
Planck Collaboration, 2020, A\&A, 641, A6,
DOI: 10.1051/0004-6361/201833910

\bibitem[Riess et al.(2022)]{riess2022}
Riess, A.G., et al., 2022, ApJL, 934, L7,
DOI: 10.3847/2041-8213/ac5c5b

\bibitem[Brout et al.(2022)]{brout2022}
Brout, D., et al., 2022, ApJ, 938, 110,
DOI: 10.3847/1538-4357/ac8e04

\bibitem[Martin(2025)]{martin2025}
Martin, E.D., 2025, Universal Horizon Address, US Patent 63/902,536,
DOI: 10.5281/zenodo.17435574

\bibitem[Riess et al.(2023)]{riess2023jwst}
Riess, A.G., et al., 2023, ApJL, 956, L18,
DOI: 10.3847/2041-8213/acf769

\bibitem[Freedman(2024)]{freedman2024}
Freedman, W.L., 2024, ApJ, 961, 32,
DOI: 10.3847/1538-4357/ad1bac

\bibitem[Poulin et al.(2019)]{poulin2019}
Poulin, V., et al., 2019, PRL, 122, 221301,
DOI: 10.1103/PhysRevLett.122.221301

\end{thebibliography}

\end{document}
